\documentclass[11pt,a4paper]{article}

\usepackage[margin=2.3cm]{geometry}
\usepackage{amsmath,amssymb}
\usepackage{booktabs}
\usepackage{xcolor}
\usepackage[most]{tcolorbox}
\usepackage[colorlinks=true,linkcolor=blue!55!black,citecolor=blue!55!black,
            urlcolor=blue!55!black]{hyperref}

\newcommand{\dalpha}{\Delta\alpha}
\newcommand{\ee}{\mathrm{e}}
\newcommand{\dd}{\mathrm{d}}
\newcommand{\ii}{\mathrm{i}}
\newcommand{\bal}{\boldsymbol{\alpha}}
\newcommand{\bphi}{\boldsymbol{\varphi}}
\newcommand{\unit}[1]{\,\mathrm{#1}}
\newcommand{\Gsf}{\mathsf{G}}

\title{\bfseries Generalised Challenge 20\\[2pt]
\large Two arbitrarily shaped particles, the multipole hierarchy, and why the
two-ellipsoid answer survives}
\author{}
\date{}

\begin{document}
\maketitle
\vspace{-2.2\baselineskip}

\begin{center}\small
Companion code: \texttt{generalized.py} \quad$\cdot$\quad
reference solution: \texttt{../answer.py}
\end{center}

\begin{tcolorbox}[colback=green!4,colframe=green!35!black,title=Short answer]
Yes.  Solve the general problem once --- two rigid dielectric bodies of
\emph{arbitrary} shape in the same tweezer geometry, each described by its full
multipole (T-matrix) response --- and the two-ellipsoid result drops out as the
leading term of a controlled double expansion in $a/\lambda$ (multipole order of
the single-particle response) and $a/R$ (multipole order of the interaction).
The reduction is sharper than a mere limit:
\begin{enumerate}\itemsep1pt
\item At dipole order the physics depends on the shape \emph{only} through the
symmetric rank-2 polarisability tensor $\bal$ and the inertia tensor
$\mathsf I$.  Since every admissible $\bal$ is realised by exactly one
ellipsoid (Sec.~\ref{sec:equiv}, closed form), \textbf{an arbitrary
sub-wavelength particle is optically indistinguishable from its ``equivalent
ellipsoid''.}  Only $\mathsf I$ has to be carried over separately.
\item What the general problem \emph{adds} is structure the challenge's
Hamiltonian hides: two librational modes per particle instead of one (four
modes, generally non-degenerate), a coupling \emph{matrix} that samples the
longitudinal Green function $G_{\parallel}\sim1/R^{2}$ as well as the transverse
$G_{\perp}\sim1/R$, an exact zero mode (spin about the polarisation), and --- for
shapes without inversion symmetry --- an $O(a/R)$ dipole--quadrupole correction
that \emph{no} ellipsoid can reproduce.
\item Set the particles axially symmetric, put the polarisation along
$\mathbf R$, and keep one libration plane: the four-mode problem collapses to
$H=\hbar\omega_t\sum_i a_i^{\dagger}a_i+\hbar g(a_1^{\dagger}a_2+\mathrm{h.c.})$
with exactly the $\omega_t$ and $g$ of \texttt{answer.py}.  Verified numerically
to $9\times10^{-16}$ relative (Sec.~\ref{sec:verify}).
\end{enumerate}
\end{tcolorbox}

\section{The general problem}

Two rigid dielectric bodies of arbitrary shape, permittivity $\epsilon_r$ and
density $\rho$, sit at positions $\mathbf r_{1,2}$ with orientations
$\Omega_{1,2}\in SO(3)$, separated by $\mathbf R=R\hat R$, each trapped by a
linearly polarised Gaussian tweezer of amplitude $E_0$, polarisation $\hat e$
and wave vector $k$.  The exact conservative light--matter energy is obtained by
solving the two-body scattering problem self-consistently; writing the response
of body $i$ as its T-matrix $\mathsf T_i(\Omega_i)$ and the free-space
propagator as $\Gsf(\mathbf R)$,
\begin{equation}
  U=-\tfrac12\operatorname{Re}\Bigl[\mathbf E_0^{\dagger}\mathsf T_1\mathbf E_0
   +\mathbf E_0^{\dagger}\mathsf T_2\mathbf E_0
   +2\,\mathbf E_0^{\dagger}\mathsf T_1\Gsf\mathsf T_2\mathbf E_0
   +O(\mathsf T^{3})\Bigr],
  \label{eq:Tmatrix}
\end{equation}
the third term being the optical binding energy at lowest order in multiple
scattering.  In the quasi-static (Rayleigh) regime it is more convenient to
expand $\mathsf T$ in Cartesian multipole polarisabilities,
\begin{equation}
  \mathbf p=\bal\cdot\mathbf E+\tfrac13\mathsf A\!:\!\nabla\mathbf E+\dots,
  \qquad
  \mathsf Q=\mathsf A\cdot\mathbf E+\mathsf C\!:\!\nabla\mathbf E+\dots,
  \label{eq:multipoles}
\end{equation}
with $\bal$ rank 2, $\mathsf A$ rank 3 and $\mathsf C$ rank 4, and to write the
interaction as
\begin{equation}
  U_{\rm int}=-C\operatorname{Re}\Bigl[\mathbf p_1\!\cdot\!\mathbf E_2
  +\tfrac13\mathsf Q_1\!:\!\nabla\mathbf E_2+\dots\Bigr],
  \qquad C=\tfrac12\ \text{(cycle average).}
  \label{eq:Uexpand}
\end{equation}
\emph{(The value of $C$ is the factor-of-two issue flagged in
\texttt{../challenge20.pdf}; it propagates unchanged into everything below.)}

\paragraph{Power counting.}
With $a$ the particle size, $\mathbf p\sim\epsilon_0a^{3}E$,
$\nabla\mathbf E_2\sim E_2/R$ and $\mathsf Q\sim\epsilon_0a^{4}E$ (if
$\mathsf A\neq0$) or $\epsilon_0a^{5}\nabla E$ (if $\mathsf A=0$).  Hence
\begin{equation}
  \frac{U_{\rm int}}{U_{\rm int}^{\rm dip-dip}}
  =1+\underbrace{O\!\left(\frac{a}{R}\right)}_{\text{only if }\mathsf A\neq0}
  +O\!\left(\frac{a}{R}\right)^{2}+\dots,
  \qquad
  \frac{\bal_{\rm exact}}{\bal_{\rm static}}=1+O\!\left(\frac{a}{\lambda}\right)^{2}.
  \label{eq:counting}
\end{equation}
\textbf{$\mathsf A$ vanishes identically for any centrosymmetric body}, an
ellipsoid included, because $\mathsf A$ is odd under $\mathbf r\to-\mathbf r$.
So for ellipsoids the series in $a/R$ is even --- which is why Appendix~A of
Ref.~\cite{Zhang2025} finds a leading correction $\propto(a^{2}-b^{2})/R^{2}$
even though it calls it ``first order''.  \emph{The odd term is the one genuinely
new effect an arbitrary shape brings to the interaction.}

\section{Dipole order in full generality}

\subsection{The equivalent ellipsoid}
\label{sec:equiv}

At leading order every sub-wavelength body is a point dipole with a symmetric
positive-definite $\bal$; let $\alpha_1\ge\alpha_2\ge\alpha_3$ be its principal
values.  For an ellipsoid with semi-axes $(a,b,c)$,
\begin{equation}
  \alpha_i=\frac{\epsilon_0V(\epsilon_r-1)}{1+L_i(\epsilon_r-1)},\quad
  V=\tfrac{4\pi}{3}abc,\quad
  L_i=\frac{abc}{2}\!\int_0^{\infty}\!\!
  \frac{\dd s}{(s+a_i^{2})\sqrt{(s+a^{2})(s+b^{2})(s+c^{2})}},
  \label{eq:Lgen}
\end{equation}
with $\sum_iL_i=1$ \cite{Osborn,Landau}.  Inverting Eq.~\eqref{eq:Lgen} and
imposing the sum rule fixes the volume in \emph{closed form},
\begin{equation}
  \boxed{\;V=\frac{\epsilon_r+2}{\epsilon_0(\epsilon_r-1)\sum_i\alpha_i^{-1}}\;},
  \qquad
  L_i=\frac{1}{\epsilon_r-1}\left[\frac{\epsilon_0V(\epsilon_r-1)}{\alpha_i}-1\right],
  \label{eq:equiv}
\end{equation}
after which the two independent axis ratios follow from the $L_i$ by a
two-dimensional root find (the map from shape to the open simplex
$\{L_i>0,\sum L_i=1\}$ is a bijection, so the ellipsoid is unique up to
relabelling).  \texttt{generalized.py} implements this and round-trips it to
$10^{-11}$.

\begin{tcolorbox}[colback=blue!4,colframe=blue!45!black]
\textbf{Equivalent-ellipsoid statement.}  At dipole order the two-particle
librational problem depends on the shape only through $\bal$ and $\mathsf I$.
Equation~\eqref{eq:equiv} produces the unique ellipsoid with the same $\bal$;
therefore \emph{the optical half of the arbitrary-shape problem is exactly the
ellipsoid problem}.  The mechanical half is not: the equivalent ellipsoid
generally has the wrong $\mathsf I$, and for a low-symmetry body the principal
axes of $\bal$ and of $\mathsf I$ need not even coincide.
\end{tcolorbox}

\subsection{Equilibrium, the zero mode, and two librations per particle}

The orientational potential of one particle in a uniform linearly polarised
field is $U(\Omega)=-\tfrac14E_0^{2}\,\hat e\cdot\bal(\Omega)\cdot\hat e$.  It is
stationary when $\hat e$ is an eigenvector of $\bal$ and minimal when $\hat e$
lies along the most polarisable body axis $\hat u_1$.  Parametrising a small
rotation by the vector $\bphi$ about equilibrium,
\begin{equation}
  U=\mathrm{const}+\tfrac14E_0^{2}
  \left[(\alpha_1-\alpha_2)\varphi_3^{2}+(\alpha_1-\alpha_3)\varphi_2^{2}\right],
  \qquad
  k_2=\tfrac12E_0^{2}(\alpha_1-\alpha_3),\;
  k_3=\tfrac12E_0^{2}(\alpha_1-\alpha_2).
  \label{eq:stiff}
\end{equation}

\begin{tcolorbox}[colback=orange!6,colframe=orange!55!black,
  title=Exact zero mode --- true for \emph{any} shape]
$\varphi_1$, the spin about the polarisation direction, does not appear in
Eq.~\eqref{eq:stiff}: rotating the body about $\hat e$ leaves
$\hat e\cdot\bal\cdot\hat e$ invariant.  A uniform, strictly linearly polarised
field therefore confines only \emph{two} of the three orientational degrees of
freedom, no matter how anisotropic the particle.  Confining the third requires
the ellipticity and longitudinal field components of a tightly focused beam, a
field gradient acting on $\mathsf C$, or the neighbour's field --- all of which
are higher order in the expansion \eqref{eq:counting}.
\end{tcolorbox}

\noindent
The two librational frequencies are
$\omega_{i\mu}=\sqrt{k_{i\mu}/I_{i\mu}}$, degenerate only when
$\alpha_2=\alpha_3$ and $I_2=I_3$, i.e.\ for an axially symmetric particle.

\subsection{The coupling matrix}

To first order in the tilts the induced dipole in the lab frame is
\begin{equation}
  \mathbf p_i=E_0\Bigl[\alpha_1\hat u_1
  +(\alpha_1-\alpha_2)\varphi_3\,\hat u_2
  -(\alpha_1-\alpha_3)\varphi_2\,\hat u_3\Bigr]+O(\varphi^{2}),
  \label{eq:dpdphi}
\end{equation}
so that
\begin{equation}
  \frac{\partial\mathbf p}{\partial\varphi_3}=(\alpha_1-\alpha_2)E_0\hat u_2,
  \qquad
  \frac{\partial\mathbf p}{\partial\varphi_2}=-(\alpha_1-\alpha_3)E_0\hat u_3 .
  \label{eq:dpdphi2}
\end{equation}
Inserting into Eq.~\eqref{eq:Uexpand} with $\mathbf E_2=\Gsf(\mathbf R)\mathbf p_2$
gives the whole dipole-order coupling in one line:
\begin{gather}
  \boxed{\;
  \kappa_{\mu\nu}=-C\operatorname{Re}
  \left[\frac{\partial\mathbf p_1}{\partial\varphi^{\mu}}\cdot
  \Gsf(\mathbf R)\cdot
  \frac{\partial\mathbf p_2}{\partial\varphi^{\nu}}\right]\;},
  \label{eq:kappa}\\[4pt]
  \Gsf(\mathbf R)=\frac{\ee^{\ii kR}}{4\pi\epsilon_0R}
  \left[\left(\mathbb 1-\hat R\hat R\right)k^{2}
  +\left(3\hat R\hat R-\mathbb 1\right)
  \left(\frac{1}{R^{2}}-\frac{\ii k}{R}\right)\right].
  \label{eq:Gdyad}
\end{gather}
The eigenvalues of $\Gsf$ are
\begin{align}
  G_{\parallel}&=\frac{2\ee^{\ii kR}(1-\ii kR)}{4\pi\epsilon_0R^{3}}
  &&\xrightarrow{\ kR\gg1\ }&&\frac{2k\sin kR}{4\pi\epsilon_0R^{2}},
  \label{eq:Gpar}\\
  G_{\perp}&=\frac{\ee^{\ii kR}(k^{2}R^{2}+\ii kR-1)}{4\pi\epsilon_0R^{3}}
  &&\xrightarrow{\ kR\gg1\ }&&\frac{k^{2}\cos kR}{4\pi\epsilon_0R},
  \label{eq:Gperp}
\end{align}
so the longitudinal channel is short ranged ($1/R^{2}$) and the transverse one
long ranged ($1/R$).

\subsection{Second quantisation}

With $\varphi_{i\mu}=\varphi_{\rm zpf}^{i\mu}(b_{i\mu}+b_{i\mu}^{\dagger})$,
$\varphi_{\rm zpf}^{i\mu}=\sqrt{\hbar/2I_{i\mu}\omega_{i\mu}}$, and the RWA,
\begin{equation}
  H=\sum_{i=1,2}\sum_{\mu=2,3}\hbar\omega_{i\mu}b_{i\mu}^{\dagger}b_{i\mu}
   +\sum_{\mu\nu}\hbar g_{\mu\nu}
    \left(b_{1\mu}^{\dagger}b_{2\nu}+\mathrm{h.c.}\right),
  \qquad
  g_{\mu\nu}=\frac{\kappa_{\mu\nu}}
  {2\sqrt{I_{1\mu}I_{2\nu}\,\omega_{1\mu}\omega_{2\nu}}} .
  \label{eq:Hgen}
\end{equation}
This is the general answer.  Everything specific to the challenge is now a
matter of specialising $\bal$, $\mathsf I$, $\hat e$ and $\hat R$.

\section{The reduction chain}

\begin{table}[h]
\centering\small
\begin{tabular}{llll}
\toprule
Step & Assumption & What collapses & Result \\
\midrule
0 & arbitrary shape, full $\mathsf T$
  & --- & Eq.~\eqref{eq:Tmatrix} \\
1 & $a\ll\lambda$, $a\ll R$
  & $\mathsf A,\mathsf C,\dots\to0$ & dipole order, Eq.~\eqref{eq:Hgen} \\
2 & optical response only
  & shape $\to\bal$ & equivalent ellipsoid, Eq.~\eqref{eq:equiv} \\
3 & $\alpha_2=\alpha_3$, $I_2=I_3$ (axial symmetry)
  & $\omega_{i2}=\omega_{i3}\equiv\omega_t$ & degenerate librations \\
4 & prolate spheroid
  & $L_{\perp}=(1-L_{\parallel})/2$ & $\dalpha$, $I=m(a^{2}+b^{2})/5$ \\
5 & $\hat e\parallel\hat R$ (the challenge)
  & $\kappa\propto\mathbb 1_{2}\,\operatorname{Re}G_{\perp}$
  & two identical $2\times2$ blocks \\
6 & keep one libration plane
  & $4\to2$ modes & $H$ of the challenge \\
\bottomrule
\end{tabular}
\caption{From the general problem to \texttt{answer.py}.}
\label{tab:chain}
\end{table}

\noindent
Step 5 is the reason the challenge's geometry is unusually clean: with
$\hat e\parallel\hat R\parallel\hat x$, \emph{both} tilt directions
$\hat u_2,\hat u_3$ are perpendicular to $\hat R$, so
$\kappa=-C(\dalpha E_0)^{2}\operatorname{Re}G_{\perp}\,\mathbb 1_{2}$: the coupling
is isotropic in the transverse plane, the off-diagonal $\kappa_{23}$ vanishes,
and the four-mode problem is two identical copies of the two-mode problem.  In
the geometry actually used in Ref.~\cite{Zhang2025} ($\hat e\perp\hat R$) one
tilt direction lies along $\hat R$ and the two modes couple through
$G_{\parallel}$ and $G_{\perp}$ respectively --- two different couplings with
different ranges (see Sec.~\ref{sec:verify}, item 2).

\section{Numerical verification}
\label{sec:verify}

\texttt{generalized.py} builds the $4\times4$ stiffness and inertia matrices
from Eqs.~\eqref{eq:stiff} and \eqref{eq:kappa} and diagonalises them.  With the
benchmark parameters ($a=300$~nm, $b=180$~nm, $\epsilon_r=5.7$,
$\rho=3500\unit{kg\,m^{-3}}$, $w_0=500$~nm, $\lambda=1064$~nm, $R=1.06\,\mu$m,
$P_0=1$~W, $C=1$):

\begin{enumerate}\itemsep3pt
\item \textbf{Reduction.}  General framework: $\omega_t/2\pi=1.344893515$~MHz for
\emph{both} librations, $g/2\pi=-100.978493209$~kHz for both diagonal entries,
$\kappa_{23}=0$ exactly.  \texttt{answer.py}: identical to
$4\times10^{-16}$ ($\omega_t$) and $9\times10^{-16}$ ($g$).
\item \textbf{Longitudinal vs.\ transverse.}  Rotating the polarisation to
$\hat y$ (the geometry of Ref.~\cite{Zhang2025}) splits the two couplings:
$g_{\parallel}/2\pi=-4.490$~kHz (in-plane tilt, $G_{\parallel}$) versus
$g_{\perp}/2\pi=-100.978$~kHz (out-of-plane tilt, $G_{\perp}$); here
$g_{\parallel}/g_{\perp}=0.0445=\operatorname{Re}G_{\parallel}/
\operatorname{Re}G_{\perp}$, a factor of $22$ that the single-mode
Hamiltonian cannot express.
\item \textbf{Exact normal modes.}  $\omega_{\pm}=1.442342$ and $1.239810$~MHz
(each doubly degenerate), matching $\omega_t\sqrt{1\pm2g/\omega_t}$ to
$4\times10^{-9}$~rad/s.  The RWA form $\omega_t\pm g$ is off by
$g^{2}/2\omega_t=3.79$~kHz, which quantifies the error of Eq.~\eqref{eq:Hgen}
itself.
\item \textbf{Triaxial particle} with semi-axes $300,220,150$~nm: the two
librations are non-degenerate, at $1.591$ and $1.010$~MHz, and carry different
couplings, $g_{22}/2\pi=-156.5$~kHz and $g_{33}/2\pi=-49.3$~kHz.  Two
frequency-resolved beam-splitter channels between the particles, not one.
\item \textbf{Equivalent ellipsoid.}  A prescribed
$\bal=(9.00,6.40,5.00)\times10^{-31}$ maps to semi-axes
$(329.9,200.4,143.0)$~nm, which reproduce $\bal$ to five digits.
\end{enumerate}

\section{Multipole corrections, and a caution about the published ones}

For a centrosymmetric particle the leading corrections are the
$O((a/R)^{2})$ and $O((a/R)^{4})$ terms of Appendix~A of Ref.~\cite{Zhang2025},
\begin{equation}
  \frac{g^{\rm exact}}{g^{\rm dip}}=1+C_1+C_2+\dots,\qquad
  C_1=\frac{3}{5}\frac{(\epsilon_r-1)(a^{2}-b^{2})}
  {\left[1+L_{\parallel}(\epsilon_r-1)\right]^{2}R^{2}},\quad
  C_2=\frac{9}{25}\frac{(\epsilon_r-1)^{2}(a^{2}-b^{2})^{2}}
  {\left[1+L_{\parallel}(\epsilon_r-1)\right]^{4}R^{4}} .
  \label{eq:C1C2}
\end{equation}

\begin{table}[h]
\centering\small
\begin{tabular}{lcc}
\toprule
Source (at $R=\lambda$, $a=300$~nm, $b=180$~nm, $\epsilon_r=5.7$)
 & $C_1$ & $C_2$ \\
\midrule
Eq.~\eqref{eq:C1C2} evaluated here            & $0.0363$ & $0.00132$ \\
Ref.~\cite{Zhang2025}, Appendix A text        & $0.0252$ & $0.0006$ \\
Ref.~\cite{Zenodo}, \texttt{FIG\_2(b).csv} header & $0.033$ & $0.012$ \\
\bottomrule
\end{tabular}
\caption{The three published/derived values of the multipole coefficients
disagree.  All are at the percent level, so none of them --- nor their
combination --- accounts for the constant factor $1.334977$ by which the Zenodo
curve exceeds Eq.~(14) of \texttt{../challenge20.pdf}.}
\label{tab:C1C2}
\end{table}

\noindent
For a shape \emph{without} inversion symmetry the expansion also contains the
odd term generated by $\mathsf A$ in Eq.~\eqref{eq:multipoles},
\begin{equation}
  \delta U_{\rm int}^{(a/R)}
  =-\tfrac{C}{3}\operatorname{Re}
  \bigl[(\mathsf A_1\!\cdot\!\mathbf E_0):\nabla\Gsf(\mathbf R)\!\cdot\!\mathbf p_2\bigr]
  +(1\leftrightarrow2),
\end{equation}
which is linear in $a/R$ and therefore \emph{dominates} $C_1$ for a strongly
asymmetric particle at close range.  This term, and a possible misalignment of
the principal axes of $\bal$ and $\mathsf I$, are the only two places where an
arbitrary shape cannot be replaced by its equivalent ellipsoid.  Both are
outside the scope of the challenge's Hamiltonian; \texttt{generalized.py}
implements the dipole order exactly and Eq.~\eqref{eq:C1C2} for reference, and
does \emph{not} implement $\mathsf A$.

\begin{thebibliography}{9}
\bibitem{Zhang2025} G.~Zhang, H.~Zhang and Z.~Yin,
Phys.\ Rev.\ A \textbf{112}, 032611 (2025),
\href{https://arxiv.org/abs/2504.08194}{arXiv:2504.08194}.
\bibitem{Zenodo} G.~Zhang, H.~Zhang and Z.~Yin (dataset),
\href{https://doi.org/10.5281/zenodo.16917848}{10.5281/zenodo.16917848} (2025).
\bibitem{Hoang2016} T.~M.~Hoang \emph{et al.},
Phys.\ Rev.\ Lett.\ \textbf{117}, 123604 (2016).
\bibitem{Rieser2022} J.~Rieser \emph{et al.}, Science \textbf{377}, 987 (2022).
\bibitem{Rudolph2024} H.~Rudolph, U.~Deli\'c, K.~Hornberger and B.~A.~Stickler,
Phys.\ Rev.\ A \textbf{110}, 063507 (2024).
\bibitem{Stickler2021} B.~A.~Stickler, K.~Hornberger and M.~S.~Kim,
Nat.\ Rev.\ Phys.\ \textbf{3}, 589 (2021).
\bibitem{Osborn} J.~A.~Osborn, Phys.\ Rev.\ \textbf{67}, 351 (1945).
\bibitem{Landau} L.~D.~Landau and E.~M.~Lifshitz,
\emph{Electrodynamics of Continuous Media}, 2nd ed., Sec.~4.
\bibitem{BH} C.~F.~Bohren and D.~R.~Huffman,
\emph{Absorption and Scattering of Light by Small Particles} (Wiley, 1983).
\end{thebibliography}

\end{document}
